\section{Proof, Not Promises}

``Formally verified'' can sound like an academic flourish. For a wealth platform
it is the opposite --- it is the most practical guarantee a buyer can get. Here is
what it actually means for your firm.

\begin{itemize}
  \item \textbf{The ledger cannot silently go out of balance.} Ratio's core
  theorem says: assemble a ledger from valid transactions, and the ledger is
  balanced --- proved for every possible sequence of inputs, by a proof assistant.
  This is not ``we tested the common cases.'' It is a guarantee that the failure
  mode every accounting system fears has been ruled out mathematically.
  \item \textbf{The engine you run is the engine that was proved.} Ratio generates
  its production code from the verified source, with a build gate that fails if the
  two ever diverge. There is no gap between ``the model we reasoned about'' and
  ``the software that shipped'' --- the classic place where verified systems quietly
  rot.
  \item \textbf{Exactness, not tolerance.} Money is integer minor units, never
  floating point. Sums are associative and exact; there is no rounding epsilon to
  tune and no end-of-day reconciliation pass to absorb drift. The number is the
  number.
  \item \textbf{Reproducibility you can hand to an auditor.} Because inputs are
  immutable and the rules are pinned by content hash, any prior state replays
  bit-for-bit. ``Show me exactly how this figure was computed, under the policy in
  force that quarter'' is a deterministic re-run, not an investigation.
\end{itemize}

\begin{callout}
\sffamily\color{ratioInk}
\textbf{The buyer's translation.} Fewer reconciliation FTEs. Faster, cleaner
audits. Restatements that are provably correct. A platform whose central numbers
you can defend to a regulator, a client, or a court --- because they are backed by
a proof, not a brand.
\end{callout}
